\section*{第 \ref{ch:g5:series-parallel-circuits} 章 --- 串联和并联电路}

\begin{solution}{exo:g5:series-parallel-circuits:1}
支路是电路分出的各条独立的
路径之一；节点是支路分开或汇合的地方。两个灯泡
并联时，各自位于一条带灯泡的支路上。
\end{solution}

\begin{solution}{exo:g5:series-parallel-circuits:2}
独立性（断开只影响自己的支路）；亮度充足
（每个灯泡都和单独接入时一样亮，电池因此耗电更快）；控制点
（支路开关只控制本支路，
分岔前的开关控制全部）。
\end{solution}

\begin{solution}{exo:g5:series-parallel-circuits:3}
灯泡 2 不受影响，仍然亮着——它自己的路径是通的。串联时
它却熄灭了，因为两个灯泡共用同一条路，
任何一处断开都会同时使两个灯泡断路。
\end{solution}

\begin{solution}{exo:g5:series-parallel-circuits:4}
B 组（并联）每个灯泡都更亮——各自达到
满亮度。B 组电池耗电也更快：
它同时为两条支路供电。
\end{solution}

\begin{solution}{exo:g5:series-parallel-circuits:5}
两个开关都闭合：两个灯泡都亮。干路开关断开：两个灯泡都灭，
因为它控制的是公共路段。干路开关闭合、支路开关断开：灯泡 1 亮，
灯泡 2 灭。
\end{solution}

\begin{solution}{exo:g5:series-parallel-circuits:6}
若是串联，一个坏掉的灯泡（或一盏被关掉的灯！）就会
使整座房子变暗，而且灯越多，所有灯就越暗。
并联接线带来的便利：每个房间的开关只控制
本房间；厨房灯泡烧坏，其余房间的灯仍然亮着；
每盏灯都能保持满亮度。
\end{solution}

\begin{solution}{exo:g5:series-parallel-circuits:7}
位于分岔之前——在电流进入住宅的公共路段上，
即保险丝盒中的总开关。它能同时切断整座房子的电路，
用于检修和紧急情况。
\end{solution}

\begin{solution}{exo:g5:series-parallel-circuits:8}
并联时：关掉天花板灯不会熄灭床头灯，
拔掉床头灯的插头也不会使天花板灯变暗——
两者各自位于自己的支路上。
\end{solution}

\begin{solution}{exo:g5:series-parallel-circuits:9}
用我们普通的开关，这些设计总会失败：与灯泡串联的开关只能从一端断开；
两个普通开关串联时，要求\emph{两个都}闭合
（任一端都能关灯，却不能总是开灯）；两个开关并联时，要使灯灭则要求
两个都断开。走廊的便利需要一种\emph{能在两条路之间做选择}的开关，
而不只是断开一条路——这就是
双控开关，
走廊中成对使用的正是它。
\end{solution}

\begin{solution}{exo:g5:series-parallel-circuits:10}
好处是：灯泡位于并联支路上，一个坏掉的灯泡
只会使它自己变暗。警告是：每少一个灯泡就少一条支路，
电池（或小彩灯串自带的小电源）会把全部输出
都送到剩下的支路上。灯泡太少时，留下的灯泡各自承受的
电流超过设计值，发热更厉害，
因而很快就烧坏了。
\end{solution}

\begin{solution}{exo:g5:series-parallel-circuits:11}
应接在支路开关自己的那段路上：位于支路开关
与灯泡 2 之间（即在该支路与其余部分汇合之前）。它
恰好在该支路工作时才亮，并且与灯泡 2 串联，
两者位于同一条支路上。
\end{solution}
